\TieudeT{PHONG TỎA VẬT LÍ}
\subsubsection*{PHẦN I. Thí sinh trả lời câu hỏi từ câu 1 đến câu 18. Mỗi câu hỏi thí sinh chỉ chọn một phương án}
\setcounter{ex}{0}
\begin{ex}%book
Một vật chuyển động, có quãng đường đi được trong khoảng thời gian t là $s$. Khi đó, tốc độ trung bình được tính bằng công thức
\choice
{$v_{t b}=\dfrac{s}{t^2}$}
{$v_{t b}=\dfrac{t}{s}$}
{\True $v_{t b}=\dfrac{s}{t}$}
{$v_{t b}=s. t$}

%\loigiai{
%}
\haidong{7} \end{ex}
\begin{ex}%book
Một người chuyển động thẳng có độ dịch chuyển $d_1$ tại thời điểm $t_1$ và độ dịch chuyển $d_2$ tại thời điểm $t_2$. Vận tốc trung bình của vật trong khoảng thời gian từ $t_1$ đến $t_2$ là
\choice
{$v=\dfrac{d_1-d_2}{t_1+t_2}$}
{\True $v=\dfrac{d_2-d_1}{t_2-t_1}$}
{$v=\dfrac{d_1+d_2}{t_2-t_1}$}
{$v=\dfrac{1}{2}\left(\dfrac{d_1}{t_1}+\dfrac{d_2}{t_2}\right)$}

%\loigiai{
%}
\haidong{7} \end{ex}
\begin{ex}%book
Đơn vị nào sau đây \textbf{không} phải đơn vị đo tốc độ?
\choice
{$km / h$}
{$m/s$}
{$km/$phút}
{\True m}

\loigiai{

}
\haidong{7} \end{ex}
\begin{ex}%book
Vào lúc $10$ giờ, người lái xe nhìn vào tốc kế và thấy tốc kế chỉ $40\ km/h$. Số liệu này cho biết
\choice
{\True tốc độ tức thời của xe}
{vận tốc trung bình của xe}
{tốc độ trung bình của xe}
{vận tốc tức thời của xe}
\loigiai{
}
\haidong{7} \end{ex}
\begin{ex}%book
Vectơ vận tốc tức thời
\choice
{\True luôn cùng hướng chuyển động}
{có gốc tại vị trí ban đầu}
{cùng chiều dương của trục $Ox$}
{có độ dài bằng độ dịch chuyển}

\loigiai{
}
\haidong{7} \end{ex}
\begin{ex}%book
Phát biểu nào sau đây về vận tốc là \textbf{sai}?
\choice
{Vận tốc là đại lượng vectơ}
{Vận tốc có thể cho biết chiều chuyển động}
{Vận tốc cho biết sự nhanh hay chậm của chuyển động}
{\True Vận tốc luôn có giá trị dương}
\loigiai{

}
\haidong{7} \end{ex}
\begin{ex}%book
Khi vật chuyển động thẳng, đổi chiều thì độ lớn của vận tốc so với tốc độ là
\choice
{bằng nhau}
{lớn hơn}
{\True nhỏ hơn}
{lớn hơn hoặc bằng}
\loigiai{

}
\haidong{7} \end{ex}
\begin{ex}%book
Chuyển động thẳng đều có tốc độ
\choice
{\True như nhau trên mọi quãng đường}
{tăng dần khi chiều dài đoạn đường tăng}
{giảm dần khi chiều dài đoạn đường tăng}
{có thể tăng dần hoặc giảm dần khi chiều dài đoạn đường tăng}
\loigiai{
}
\haidong{7} \end{ex}
\begin{ex}%Tailieuchuan
Một ô tô chuyển động trên đường thẳng và không đổi chiều. Tại thời điểm $t_1$, ô tô ở cách vị trí xuất phát $8\ km$. Tại thời điểm $t_2$, ô tô cách vị trí xuất phát $15\ km$. Từ $t_1$ đến $t_2$, độ dịch chuyển của ô tô đã thay đổi một đoạn bằng
\choice
{$24\ km$}
{$8\ km$}
{\True $7\ km$}
{$15\ km$}
\loigiai{

}
\end{ex}

\begin{ex}%book
Trong thí nghiệm đo tốc độ trung bình của viên bi thép chuyển động trên máng nghiêng dài $0,5\ m$, người ta tính được giá trị trung bình của thời gian chuyển động trong $3$ lần đo là $0,778\ s$. Tốc độ trung bình của viên bi có giá trị là
\choice
{\True $0,643\ m/s$}
{$0,625\ m/s$}
{$0,647\ m/s$}
{$0,629\ m/s$}
\haidong{12}
\loigiai{
}
\end{ex}
\begin{ex}%book%[1P1K2-3][Loai1] 
Một xe máy đi nửa đoạn đường đầu tiên với tốc độ trung bình $v_1 = 60\ km/h$ và nửa đoạn đường sau với tốc độ trung bình $v_2 = 40\ km/h$. Tốc độ trung bình trên cả đoạn đường là
\choice
{$35\ km/h$}
{$44\ km/h$}
{$65\ km/h$}
{\True $48\ km/h$}
\haidong{12}
\loigiai{
\begin{itemize}
\item Gọi toàn bộ quãng đường là 2 s. 
\item Nửa quãng đường đầu là $s_1$ và nửa quãng đường sau là $s_2$. 
\item Theo đề ra ta có: $s_1=s_2=s$.
\item Thời gian đi trên nửa đường đầu và nửa đường sau là: $t_1=\dfrac{s}{v_1}$ và $t_2=\dfrac{s}{v_2}$.
\item Tốc độ trung bình trên toàn bộ quãng đường:
\begin{align*}
v=\dfrac{s_1+s_2}{t_1+t_2}=\dfrac{2s}{\dfrac{s}{v_1}+\dfrac{s}{v_1}}=\dfrac{2v_1v_2}{v_1+v_2}=\dfrac{2.60.40}{60+40}=48\ km/h.
\end{align*}
\end{itemize}}

%<MyLT>
\end{ex}
\begin{ex}
Cho quãng đường $s=0,5\ m$ và bảng kết quả thí nghiệm. Tốc độ trung bình là
\begin{center}
\begin{tabular}{|c|c|c|c|c|}
\hline
Lần đo & Lần $1$ & Lần $2$ & Lần $3$ & Giá trị trung bình \\
\hline
Thời gian $t\ (s)$ & $0,777$ & $0,780$ & $0,776$ &  \\
\hline
\end{tabular}
\end{center}

\choice
{\True $0,643\ m/s$}
{$0,732\ m/s$}
{$0,824\ m/s$}
{$0,657\ m/s$}
\haidong{10}
\loigiai{
}
\end{ex}
\begin{ex}
Để xác định tốc độ của một vật chuyển động đều, một người đã đo quãng đường vật đi được bằng $(16,0 \pm 0,4)\ m$ trong khoảng thời gian là $(4,0 \pm 0,2)\ s$. Tốc độ của vật là
\choice
{\True $(4,0 \pm 0,3)\ m/s$}
{$(4,0 \pm 0,6)\ m/s$}
{$(4,0 \pm 0,2)\ m/s$}
{$(4,0 \pm 0,1)\ m/s$}
\haidong{13}
\loigiai{

\begin{itemize}
\item ta có $v=\dfrac{s}{t} \Rightarrow \overline{v}=\dfrac{\overline{s}}{t}=\dfrac{16}{4}=4\ m/s$

\item $\dfrac{\Delta v}{\overline{v}}=\dfrac{\Delta s}{\overline{s}}+\dfrac{\Delta t}{\overline{t}}=\dfrac{0,4}{16}+\dfrac{0,2}{4}=0,075 \Rightarrow \Delta v=0,075, \overline{v}=0,3\ m/s \Rightarrow v=(4,0 \pm 0,3) m/s$

\end{itemize}


}
\end{ex}
\begin{ex}%book
Trong thí nghiệm thực hành đo tốc độ của vật chuyển động, sử dụng hai cổng quang điện để đo
\choice
{\True thời gian chuyển động của viên bi thép}
{tốc độ trung bình của viên bi thép}
{đường kính của viên bi thép}
{tốc độ tức thời của viên bi thép}

\loigiai{
}
\haidong{7} \end{ex}
\begin{ex}%book
Sử dụng đồng hồ đo thời gian hiện số và cổng quang điện để đo tốc độ chuyển động phát biểu nào sau đây là \textbf{sai}?
\choice
{Ưu điểm: Độ chính xác cao đến hàng nghìn giây}
{Nhược điểm: Thiết bị đo cồng kềnh}
{Nhược điểm: Tốn kém kinh phí}
{\True Ưu điểm: Thiết bị đo gọn nhẹ}
\loigiai{

}
\haidong{7} \end{ex}
\begin{ex}%book
Dụng cụ nào sau đây \textbf{không} dùng để thực hiện thí nghiệm đo tốc độ trong phòng thí nghiệm?
\choice
{Đồng hồ đo thời gian hiện số}
{Cổng quang điện}
{\True Ampe kế}
{Viên bi thép}

\loigiai{
}
\haidong{7} \end{ex}
\begin{ex}%book
Dụng cụ nào sau đây \textbf{không} có trong thực hành đo tốc độ chuyển động của một viên bi sử dụng cổng quang điện?
\choice
{Viên bi thép}
{\True Thước kẹp để đo đường kính viên bi}
{Cổng quang điện}
{Đồng hồ bấm giây}
\loigiai{
}
\haidong{7} \end{ex}
\begin{ex}%book
Để xác định tốc độ của một vật chuyển động thẳng đều, một nhóm học sinh tiến hành đo quãng đường s vật đi được trong khoảng thời gian t. Kết quả đo được quãng đường $s=(24,0 \pm 0,2)\ m$ và thời gian $t=(4,0 \pm 0,1)\ s$. Tốc độ của vật là
\choice
{$v=(4,0 \pm 0,2)\ m/s$}
{$v=(6,0 \pm 0,3)\ m/s$}
{$v=(4,0 \pm 0,3)\ m/s$}
{\True $v=(6,0 \pm 0,2)\ m/s$}
\haidong{13}
\loigiai{
}
\haidong{7} \end{ex}



































\subsubsection*{PHẦN 2. Thí sinh trả lời từ câu 1 đến câu 4. Trong mỗi ý a), b), c), d) ở mỗi câu, thí sinh chọn đúng hoặc sai.}
\setcounter{ex}{0}
\begin{ex}%Tailieuchuan
Một ô tô chạy từ địa điểm $A$ đến địa điểm $B$ với tốc độ $40 \ km/h$, sau đó trở về $A$ với tốc độ $60\ km/h$.
\choiceTF[t]
{\True Ô tô đi từ $A$ đến $B$ chậm hơn đi từ $B$ trở về $A$} 
{\True Thời gian ô tô đi từ $A$ đến $B$ gấp $1,5$ lần thời gian ô tô đi từ $B$ về $A$} 
{Tốc độ trung bình trên cả quãng đường (đi và về) của ô tô là $50 \ km/h$} 
{\True Vận tốc trung bình trên cả quãng đường (đi và về) của ô tô là $0$} 
\loigiai{ 
\begin{itemchoice}
\itemch
\itemch
\itemch
\itemch
\end{itemchoice}
}
\end{ex}
\begin{ex}%Tailieuchuan
Nhân dịp nghỉ lễ Quốc tế lao động, gia đình nhà bạn Tùng quyết định về quê nghỉ lễ. Biết gia đình nhà Tùng di chuyển bằng xe ô tô chuyển động đều từ Hà Nội về quê Thái Bình với quãng đường $130\ km$. Xe ô tô xuất phát từ Hà Nội lúc $6$ giờ sáng và chạy liên tục không nghỉ với tốc độ giả sử không đổi là $13\ m/s$ trên suốt quãng đường.
\choiceTF[t]
{\True Độ dịch chuyển của xe luôn tăng dần theo thời gian} 
{\True Tốc độ tức thời của xe tại thời điểm $7$ giờ sáng là $13 \ m/s$} 
{Tốc độ trung bình của xe trên cả quãng đường là $13 \ km/h$} 
{Sau $2,5$ giờ thì gia đình nhà bạn Tùng về tới quê Thái Bình} 
\loigiai{ 
\begin{itemchoice}
\itemch
\itemch
\itemch
\itemch
\end{itemchoice}
}
\end{ex}
\begin{ex}%Tailieuchuan
Một ô tô chở khách trong hành trình $5$ giờ đi từ Hà Nội đến Bắc Ninh. Biết ô tô khởi hành lúc $6$ giờ sáng với quãng đường đi được khi tới nơi là $40,5 \ km$. Xem tốc độ của ô tô trên cả quãng đường là không đổi.
\choiceTF[t]
{\True Ô tô đến Bắc Ninh vào lúc $11$ giờ} 
{Tốc độ trung bình của ô tô trên cả quãng đường là $8,5\ km/h$} 
{\True Tốc độ trung bình của ô tô bằng độ lớn vận tốc trung bình của ô tô} 
{\True Lúc $7$ giờ có một xe máy cũng bắt đầu xuất phát từ Hà Nội với lộ trình như ô tô. Khoảng cách lúc này của hai xe lớn hơn $8\ km$} 
\loigiai{ 
\begin{itemchoice}
\itemch
\itemch
\itemch
\itemch
\end{itemchoice}
}
\end{ex}
\begin{ex}%Tailieuchuan
\immini[thm]{Bạn Hoa đi học từ nhà $A$ đến trường $C$ theo lộ trình $ABC$ như hình vẽ. Bạn Hoa đi đoạn đường $AB=400 \ m$ hết $6$ phút, đoạn đường $BC=300 \ m$ hết $4$ phút.}{\begin{tikzpicture}[scale=1,thick,>=stealth',draw=Mapcolor]
\path
(0,0)coordinate(A)
(3,0)coordinate(B)
(3,2.5)coordinate(C)
;
\draw[Mapcolor] (A)--(B) (B)--(C);
\foreach \x/\g in {A/180,B/-90,C/90}
\draw[fill=white](\x) circle(0.05)+(\g:0.3)node{$\x$};
\end{tikzpicture}}
\choiceTF[t]
{\True Khi bạn Hoa đi từ $B$ đến $C$, độ lớn của độ dịch chuyển cũng chính bằng quãng đường mà Hoa đi được} 
{Độ dịch chuyển của bạn Hoa khi đi từ nhà đến trường có độ lớn là $700\ m$} 
{Tốc độ trung bình của bạn Hoa khi đi từ nhà đến trường là $\dfrac{7}{6}$ m/phút} 
{Vận tốc trung bình của bạn Hoa khi đi từ nhà đến trường là $\dfrac{25}{12} \ m/s$} 
\loigiai{ 
\begin{itemchoice}
\itemch
\itemch
\itemch
\itemch
\end{itemchoice}
}
\end{ex}


\subsubsection*{PHẦN 3. Thí sinh trả lời từ câu 1 đến câu 6.}
\setcounter{ex}{0}
\begin{ex}%Tailieuchuan
Một máy bay phản lực bay theo đường thẳng với tốc độ trung bình là $2400\ km/h$. Nếu muốn bay liên tục trên khoảng cách $6000\ km$ thì máy bay phải bay trong bao nhiêu giờ?
\shortans[oly]{2,5}
\loigiai{
}
\end{ex}
\begin{ex}%Tailieuchuan
Một người đi xe máy chuyển động thẳng đều từ $A$ lúc $5$ giờ sáng và tới $B$ lúc $7$ giờ $30$ phút, biết độ dài quãng đường $AB$ là $150\ km$. Tới $B$ xe dừng lại $45$ phút rồi đi về $A$ với vận tốc $50\ km/h$. Xe tới $A$ lúc $11$ giờ $x$ phút. Tìm $x$.
\shortans[oly]{15}
\loigiai{
}
\end{ex}
\begin{ex}%Tailieuchuan
Một vật chuyển động thẳng không đổi chiều trên quãng đường dài $35\ m$. Nửa quãng đường đầu vật đi hết thời gian $5\ s$, nửa quãng đường sau vật đi hết thời gian $2\ s$. Tốc độ trung bình trên cả quãng đường là bao nhiêu $km/h$?
\shortans[oly]{18}
\loigiai{
}
\end{ex}
\begin{ex}%Tailieuchuan
Một xe chuyển động thẳng không đổi chiều có tốc độ trung bình là $20\ km/h$ trên $\dfrac{1}{4}$ đoạn đường đầu và $40\ km / h$ trên đoạn đường còn lại. Tốc độ trung bình của xe trên cả đoạn đường là bao nhiêu $km/h$?
\shortans[oly]{32}
\loigiai{
}
\end{ex}
\begin{ex}%Tailieuchuan
Một chiếc xe chạy $50\ km$ đầu tiên với vận tốc $25\ km/h$; $70\ km$ sau với tốc độ $35\ km/h$. Tốc độ trung bình của xe trong suốt quãng đường chuyển động là bao nhiêu $km/h$ (làm tròn kết quả đến chữ số hàng đơn vị)?
\shortans[oly]{30}
\loigiai{}
\end{ex}
\begin{ex}%Tailieuchuan
Một vật chuyển động với tốc độ $v$ đi hết quãng đường $42\ km$ mất thời gian là $t$. Nếu tốc độ tăng gấp đôi, nhưng thời gian thời gian giảm bớt $\dfrac{1}{3}$ thì vật đi được quãng đường là bao nhiêu kilomet (làm tròn kết quả đến chữ số hàng đơn vị)?
\shortans[oly]{56}
\loigiai{
}
\end{ex}






\newpage 


\TieudeT{PHONG TỎA VẬT LÍ}
\subsubsection*{PHẦN I. Thí sinh trả lời câu hỏi từ câu 1 đến câu 18. Mỗi câu hỏi thí sinh chỉ chọn một phương án}
\setcounter{ex}{0}
\begin{ex}%book
Người ta thường dùng quãng đường đi được trong cùng một đơn vị thời gian để xác định độ nhanh, chậm của chuyển động. Đại lượng này gọi là
\choice
{Vận tốc tức thời}
{Vận tốc trung bình}
{Tốc độ tức thời}
{\True Tốc độ trung bình}

%\loigiai{
%}
\haidong{7} \end{ex}
\begin{ex}%book
Vận tốc trung bình là đại lượng được đo bởi
\choice
{\True thương số của độ dịch chuyển và thời gian dịch chuyển}
{tích của quãng đường đi được và thời gian dịch chuyển}
{tích của độ dịch chuyển và thời gian dịch chuyển}
{thương số của quãng đường đi được và khoảng thời gian đi hết quãng đường}
\loigiai{

}
\haidong{7} \end{ex}
\begin{ex}%book
Tốc độ trung bình trong một thời gian rất ngắn được gọi là
\choice
{vận tốc}
{tốc độ}
{\True tốc độ tức thời}
{vận tốc trung bình}
\loigiai{

}
\haidong{7} 
\end{ex}
\begin{ex}%book
Khi vật chuyển động có độ dịch chuyển $\Delta \overrightarrow{d}$ trong khoảng thời gian $\Delta t$ rất nhỏ. Vận tốc tức thời của vật được tính bằng
\choice
{\True $\overrightarrow{v}_t=\dfrac{\Delta \overrightarrow{d}}{\Delta t}$}
{$\overrightarrow{v}_t=\Delta \overrightarrow{d}. \Delta t$}
{$\overrightarrow{v}_t=\dfrac{\Delta t}{\Delta \vec{d}}$}
{$\overrightarrow{v}_t=\Delta \overrightarrow{d}+\Delta t$}

\loigiai{
}
\haidong{7} \end{ex}
\begin{ex}%book
Trong chuyển động thẳng, vectơ vận tốc tức thời có
\choice
{phương và chiều luôn thay đổi}
{phương không đổi, chiều có thể thay đổi}
{\True phương và chiều không thay đổi}
{phương không đổi, chiều luôn thay đổi}
\loigiai{
}
\haidong{7} \end{ex}
\begin{ex}%book
Phát biểu nào sau đây về vận tốc là đúng?
\choice
{Vận tốc là đại lượng vectơ có hướng ngược hướng với hướng của độ dịch chuyển}
{Vận tốc là đại lượng vô hướng có thể âm hoặc dương}
{\True Vận tốc là đại lượng vectơ có hướng là hướng của độ dịch chuyển}
{Vận tốc là đại lượng vô hướng không âm}
\loigiai{

}
\haidong{7} \end{ex}
\begin{ex}%book
Phát biểu nào sau đây là đúng?
\choice
{\True Vận tốc trung bình là một đại lượng có hướng}
{Vận tốc trung bình là một đại lượng vô hướng}
{Tốc độ trung bình là một đại lượng có hướng}
{Tốc độ tức thời là một đại lượng có hướng}
\loigiai{
}
\haidong{7} \end{ex}
\begin{ex}%Tailieuchuan
Một vận động viên chạy theo bốn chặng như sau:
\begin{itemize}
\item Chặng $1$ dài $200\ m$ với mất $20\ s$. 
\item Chặng $2$ dài $100\ m$ mất $10\ s$.
\item Chặng $3$ dài $400\ m$ mất $43,5\ s$.
\item Chặng $4$ dài $150\ m$ mất $13\ s$.
\end{itemize}  
Vận động viên chạy nhanh nhất ở
\choice
{chặng $1$}
{chặng $2$}
{chặng $3$}
{\True chặng $4$}
\loigiai{		
}
\end{ex}
\begin{ex}%Tailieuchuan
Một người đi xe máy từ nhà đến hiệu sách với quãng đường $3\ km$ mất $15$ phút, sau đó đi đến trường với quãng đường $2\ km$ mất $12$ phút. Biết trường nằm giữa nhà và nhà sách và cùng nằm trên một đường thẳng. Vận tốc trung bình của người đi xe máy là
\choice
{\True $11,1\ km/h$}
{$0,98\ m/s$}
{$0,29\ km / h$}
{$0,38\ m/s$}
\loigiai{

}
\end{ex}
\begin{ex}%book
Một người đi xe máy đi từ ngã tư với tốc độ trung bình $30\ km/h$ theo hướng Bắc. Sau $3$ phút người đó đến vị trí nào trên hình?
\begin{center}
\begin{tikzpicture}
\node[] at (0,0){\includegraphics[width=3cm]{images/5-1}};
\end{tikzpicture}
\end{center}
\choice
{\True $E$}
{$G$}
{$I$}
{$C$}
\loigiai{
\begin{itemize}
\item Đổi: 3 phút = 0,05 giờ
\item Quãng đường người đó đi được sau 3 phút là: $ s=vt=30.0,05=1,5\ km$.
\item Vậy sau 3 phút, người đó đến vị trí E trên hình.
\end{itemize}
}
\haidong{13}
%<MyLT1>
\end{ex}
\begin{ex}%book
Một người đi xe máy từ nhà đến siêu thị mất $0,25\ h$, sau đó trở về nhà trong thời gian $0,2\ h$. Hai địa điểm cách nhau $9\ km$. Coi quỹ đạo đi được là đường thẳng. Tốc độ trung bình của người đó là
\choice
{$40,5\ km/h$}
{$20\ km/h$}
{$40\ m/s$}
{\True $40\ km/h$}
\haidong{12}
\loigiai{
}
\end{ex}
\begin{ex}%book
Một người đi xe đạp trên $\dfrac{2}{3}$ đoạn đường đầu với tốc độ trung bình $15\ km/h$ và $\dfrac{1}{3}$ đoạn đường sau với tốc độ trung bình $20\ km/h$. Tốc độ trung bình của người đi xe đạp trên cả quãng đường là
\choice
{$17,5\ km/h$}
{$12\ km/h$}
{$15\ km/h$}
{\True $16,36\ km/h$}
\haidong{12}
\loigiai{
\begin{itemize}
\item Gọi $t_1$ và $t_2$ lần lượt là thời gian vật chuyển động với vận tốc $15\ km/h$ và $20\ km/h$.
\item Ta có: $15 \cdot t_1=2 \cdot 20 \cdot t_2 \Rightarrow t_2=\dfrac{3}{8} t_1$
\item Tốc độ trung bình trên cả đoạn đường là $v=\dfrac{\Delta s}{\Delta t}=\dfrac{20 \cdot t_2+15 \cdot t_1}{t_2+t_1}=\dfrac{22,5 \cdot t_1}{\dfrac{11}{8} t_1}=16,36\ km/h$.
\end{itemize}
}
\end{ex}
\begin{ex}
Để xác định tốc độ của một vật chuyển động đều, một người đã đo quãng đường vật đi được bằng $20,0 \pm 0,4\ m$ trong khoảng thời gian là $4,0 \pm 0,2\ s$. Sai số tuyệt đối của vận tốc là
\choice
{$0,3\ m/s$}
{$0,4\ m/s$}
{\True $0,35\ m/s$}
{$0,45\ m/s$}
\haidong{13}
\loigiai{
}
\end{ex}
\begin{ex}%book
Ưu điểm của việc sử dụng đồng hồ đo thời gian hiện số để đo tốc độ chuyển động là
\choice
{\True đo thời gian chính xác đến phần nghìn giây}
{thao tác thực hiện phức tạp}
{thời gian bắt đầu đo kém chính xác}
{cấu tạo cồng kềnh, phức tạp}

\loigiai{
}
\haidong{7} \end{ex}
\begin{ex}%book
Ưu điểm khi sử dụng đồng hồ đo thời gian hiện số và cổng quang điện là
\choice
{\True kết quả có độ chính xác cao}
{thiết bị nhỏ, gọn}
{tuổi thọ cao}
{chi phí thấp}

\loigiai{
}
\haidong{7} \end{ex}
\begin{ex}%book
Sử dụng đồng hồ đo thời gian hiện số để xác định thời gian chuyển động của vật từ A đến B ta điều chỉnh núm xoay về MODE nào sau đây?
\choice
{$MODE\ A$}
{$MODE\ B$}
{\True $MODE\ A \leftrightarrow B$}
{$MODE\ A+B$}
\loigiai{
}
\haidong{7} \end{ex}
\begin{ex}%book
Để đo tốc độ tức thời và tốc độ trung bình của một vật chuyển động trong phòng thí nghiệm ta có thể dùng cổng quang điện kết hợp với
\choice
{\True thước và đồng hồ đo thời gian hiện số}
{nam châm điện, quả dọi và đồng hồ bấm giây}
{đồng hồ đo thời gian hiện số và nam châm điện}
{thước và nam châm điện}

\loigiai{
}
\haidong{7} \end{ex}
\begin{ex}%book
Muốn đo tốc độ trung bình của một vật chuyển động thẳng thì người ta cần đo những đại lượng nào?
\choice
{Đo đường kính và khối lượng của vật}
{\True Đo quãng đường và thời gian chuyển động của vật}
{Đo khối lượng và thời gian chuyển động của vật}
{Đo quãng đường và khối lượng của vật}
\loigiai{

}
\haidong{7} \end{ex}













































\subsubsection*{PHẦN 2. Thí sinh trả lời từ câu 1 đến câu 4. Trong mỗi ý a), b), c), d) ở mỗi câu, thí sinh chọn đúng hoặc sai.}
\setcounter{ex}{0}
\begin{ex}%book
Hai học sinh chở nhau bằng xe đạp điện từ trường THPT Chuyên Quốc Học dọc theo đường Lê Lợi đến quán chè Hẻm trên đường Hùng Vương như hình bên hết thời gian $5$ phút. 
\immini{\choiceTF[t]
{Độ dịch chuyển của xe là $1,5\ km$}
{\True Vận tốc trung bình của xe là $13,4\ km/h$}
{\True Quãng đường đi được của xe là $1,5\ km$}
{\True Xe chuyển động với tốc độ trung bình $18\ km/h$}
}{\includegraphics[width=4.25cm]{images/4_11}}
\loigiai{
\begin{itemchoice}
\itemch
\itemch
\itemch
\itemch
\end{itemchoice}
}
\end{ex}
\begin{ex}%Tailieuchuan
Một xe buýt chuyển động thẳng không đổi chiều trên quãng đường từ $A$ đến $B$ dài $350\ m$. Nửa quãng đường đầu vật đi hết thời gian $50$ giây, nửa quãng đường sau vật đi hết thời gian $20$ giây.
\choiceTF[t]
{\True Nếu nửa quãng đường đầu, xe buýt chạy với tốc độ bằng tốc độ trên nửa quãng đường sau thì tổng thời gian chuyển động sẽ giảm $30$ giây} 
{\True Tốc độ của vật trên nửa quãng đường đầu nhỏ hơn trên nửa quãng đường sau} 
{Tốc độ trung bình của vật trên cả quãng đường nhỏ hơn tốc độ của vật trên nửa quãng đường đầu} 
{Nếu một ô tô cũng xuất phát từ $A$ cùng lúc với xe buýt, chuyển động với tốc độ không đổi là $4,7 \ m/s$ trên cả quãng đường thì sẽ đến $B$ trước xe buýt} 
\loigiai{ 
\begin{itemchoice}
\itemch
\itemch
\itemch
\itemch
\end{itemchoice}
}
\end{ex}
\begin{ex}%Tailieuchuan
Một ô tô chạy thẳng từ địa điểm $A$ đến địa điểm $B$ với tốc độ $40 \ km / h$, sau đó quay về $A$ với tốc độ $60 \ km/h$. Giả sử tốc độ của ô tô luôn không đổi.
\choiceTF[t]
{Vận tốc là đại lượng vectơ, có thể bằng $0$ hoặc dương, không có giá trị âm} 
{Thời gian ô tô đi từ $A$ đến $B$ và thời gian đi từ $B$ trở về $A$ là bằng nhau} 
{\True Độ dịch chuyển trên cả quãng đường đi và về của ô tô là $0\ km$} 
{\True Tốc độ trung bình trên cả quãng đường đi và về của ô tô là $48\ km/h$} 
\loigiai{ 
\begin{itemchoice}
\itemch
\itemch
\itemch
\itemch
\end{itemchoice}
}
\end{ex}




\begin{ex}%Tailieuchuan
Một tàu ngầm sử dụng hệ thống phát sóng âm để đo độ sâu $h$ của biển. Hệ thống phát ra sóng âm và đo thời gian quay trở lại của sóng âm sau khi chúng bị phản xạ tại đáy biển. Tại một vị trí trên mặt biển, thời gian mà hệ thống ghi nhận được là $0,13 \ s$ kể từ khi sóng âm được truyền đi. Biết tốc độ truyền sóng âm trong nước là $1500 \ m/s$.
\choiceTF[t]
{Quãng đường sóng âm truyền bằng độ sâu của mực nước biển} 
{Độ dịch chuyển của sóng âm chính bằng quãng đường sóng âm truyền} 
{\True Tốc độ trung bình của sóng âm truyền trong nước kể từ khi sóng âm được truyền đi đến khi hệ thống ghi nhận được là $1500 \ m/s$} 
{\True Độ sâu của mực nước biên là $97,5 \ m$} 
\loigiai{ 
\begin{itemchoice}
\itemch
\itemch
\itemch
\itemch
\end{itemchoice}
}
\end{ex}


\subsubsection*{PHẦN 3. Thí sinh trả lời từ câu 1 đến câu 6.}
\setcounter{ex}{0}
\begin{ex}%Tailieuchuan
Một vật chuyển động trên quãng đường $s$ mất $3$ giờ. Nếu phải đi quãng đường dài gấp $3$ lần nhưng tốc độ trung bình chỉ tăng gấp đôi thì phải mất bao nhiêu giờ (làm tròn kết quả đến chữ số hàng phần mười)?
\shortans[oly]{4,5}
\loigiai{
}
\end{ex}
\begin{ex}%Tailieuchuan
Một người đi với tốc độ $60\ km/h$ trên nửa đoạn đường đầu, trên nửa đoạn đường còn lại, người đó đi nửa thời gian đầu với tốc độ $40\ km/h$ và nửa thời gian sau với tốc độ $20\ km/h$. Tốc độ trung bình trên cả đoạn đường là bao nhiêu $km/h$ (làm tròn kết quả đến chữ số hàng đơn vị)?
\shortans[oly]{40}
\loigiai{
}
\end{ex}
\begin{ex}%Tailieuchuan
Khi lái xe trên đường, tài xế chỉ mất tập trung một khoảng thời gian rất nhỏ có thể gây ra va chạm không mong muốn. Khi một người hắt hơi mạnh, mắt của người đó có thể nhắm mắt lại trong $0,5\ s$. Nếu người đó đang lái xe với tốc độ $90\ km/h$ thì trong thời gian nhắm mắt đó xe sẽ đi được quãng đường bao nhiêu mét?
\shortans[oly]{12,5}
\loigiai{
}
\end{ex}
\begin{ex}%Tailieuchuan
Một người đi bộ $3\ km$ trên một con đường thẳng theo hướng Bắc rồi quay đầu lại và đi $1\ km$ theo hướng Nam và tổng thời gian đi của người này là $1$ giờ. Tốc độ trung bình của người này là bao nhiêu $m/s$ (làm tròn kết quả đến chữ số hàng phần trăm)?
\shortans[oly]{1,11}
\loigiai{
}
\end{ex}
\begin{ex}%Tailieuchuan
Để xác định tốc độ của một vật chuyển động đều, một người đã đo quãng đường vật đi được bằng $(16,0 \pm 0,4)\ m$ trong khoảng thời gian là $(4,0 \pm 0,2)\ s$. Sai số phép đo tốc độ trung bình của vật là bao nhiêu $m/s$?
\shortans[oly]{0,3}
\loigiai{
}
\end{ex}
\begin{ex}%Tailieuchuan
Một ô tô xuất phát từ $A$ lúc $6$ giờ sáng, chuyển động thẳng đều tới $B$, cách $A$ $90\ km$, biết rằng xe tới $B$ lúc $8$ giờ $30$ phút. Tốc độ của xe là bao nhiêu $km/h$?
\shortans[oly]{36}
\loigiai{
}
\end{ex}




\newpage 


\TieudeT{PHONG TỎA VẬT LÍ}
\subsubsection*{PHẦN I. Thí sinh trả lời câu hỏi từ câu 1 đến câu 18. Mỗi câu hỏi thí sinh chỉ chọn một phương án}
\setcounter{ex}{0}
\begin{ex}%book
Công thức tính vận tốc trung bình của một vật chuyển động là
\choice
{$\vec{v}=\dfrac{\vec{s}}{t}$}
{\True $\vec{v}=\dfrac{\vec{d}}{t}$}
{$v=\dfrac{d}{t}$}
{$v=\dfrac{s}{t}$}

\loigiai{

}
\haidong{7} \end{ex}
\begin{ex}%book
Đơn vị đo tốc độ là
\choice
{\True $m/s$}
{m}
{$m/s^2$}
{$s / m$}

\loigiai{

}
\haidong{7} \end{ex}
\begin{ex}%book
Phát biểu nào sau đây là đúng khi nói về tốc độ tức thời?
\choice
{Tốc độ tức thời diễn tả sự nhanh chậm của chuyển động trên cả quãng đường vật đi được}
{\True Tốc độ tức thời diễn tả sự nhanh chậm của chuyển động tại một thời điểm xác định}
{Tốc độ tức thời là tốc độ trung bình trong toàn bộ thời gian chuyển động}
{Tốc độ tức thời là cách gọi khác của tốc độ trung bình}

\loigiai{
}
\haidong{7} \end{ex}
\begin{ex}%book
Vận tốc tức thời là
\choice
{\True vận tốc trung bình xét trong một khoảng thời gian rất nhỏ trong quá trình chuyển động của vật}
{tốc độ trung bình xét trong một khoảng thời gian rất nhỏ trong quá trình chuyển động của vật}
{tốc độ trung bình xét trong toàn bộ quá trình chuyển động của vật}
{vận tốc trung bình xét trong toàn bộ quá trình chuyển động của vật}
\loigiai{
}
\haidong{7} 
\end{ex}
\begin{ex}%book
Tốc độ là đại lượng đặc trưng cho
\choice
{\True tính chất nhanh hay chậm của chuyển động}
{sự thay đổi hướng của chuyển động}
{khả năng duy trì chuyển động của vật}
{sự thay đổi vị trí của vật trong không gian}
\loigiai{
}
\haidong{7} \end{ex}
\begin{ex}%book
Hai đại lượng nào sau đây là đại lượng vectơ?
\choice
{Quãng đường và độ dịch chuyển}
{\True Độ dịch chuyển và vận tốc}
{Quãng đường và tốc độ}
{Tốc độ và vận tốc}

\loigiai{
}
\haidong{7} \end{ex}
\begin{ex}%book
Bố bạn A đưa A đi học bằng xe máy, xuất phát lúc $7$ giờ đến trường trên quãng đường $15\ km$. Đến gần trường, xe giảm dần tốc độ và dừng trước cổng trường lúc $7$ giờ $30$ phút. Tốc độ trung bình của xe trên quãng đường đến trường là
\choice
{\True $30\ km/h$}
{$20\ km/h$}
{$10\ km/h$}
{$40\ km/h$}
\haidong{12}
\loigiai{
}
\end{ex}
\begin{ex}%book
Khi vật chuyển động thẳng đổi chiều, thì trong khoảng thời gian ngược chiều đó 
\choice
{quãng đường đi được vẫn bằng độ dịch chuyển}
{tốc độ có giá trị âm, vận tốc có giá trị dương}
{\True tốc độ và vận tốc có độ lớn bằng nhau}
{tốc độ có giá trị dương còn vận tốc có giá trị âm}
\loigiai{
}
\haidong{7} \end{ex}

\begin{ex}%book
Một xe chạy liên tục trong $2,5$ giờ, trong $1$ giờ đầu, tốc độ trung bình của xe là $60\ km/h$, trong $1,5$ giờ sau, tốc độ trung bình của xe là $40\ km/h$. Tốc độ trung bình của xe trong toàn bộ khoảng thời gian chuyển động là
\choice
{\True $48\ km/h$}
{$37\ km/h$}
{$54\ km/h$}
{$45\ km/h$}
\haidong{12}
\loigiai{
}
\end{ex}
\begin{ex}%Tailieuchuan
Một xe chuyển động thẳng không đổi chiều; $2$ giờ đầu xe chạy với tốc độ trung bình $40\ km/h$, $3$ giờ sau xe chạy với tốc độ trung bình $25\ km/h$. Tốc độ trung bình của xe trong suốt thời gian chạy là
\choice
{$30\ km/h$}
{\True $31\ km/h$}
{$32\ km/h$}
{$33\ km/h$}
\loigiai{}
\end{ex}

\begin{ex}%Tailieuchuan
Một vận động viên chạy một quãng đường $4\ km$ theo $3$ chặng như sau: 
\begin{itemize}
\item 	Nửa quãng đường đầu anh ta chạy hết thời gian $6$ phút $40$ giây.
\item $\dfrac{2}{3}$ quãng đường còn lại anh ta chạy hết thời gian $5$ phút $33$ giây. 
\item $\dfrac{1}{3}$ đoạn đường cuối chạy hết $1$ phút $45$ giây.
\end{itemize}
Tốc độ trung bình của vận động viên đó trên cả quãng đường xấp xỉ
\choice
{$9,5\ m/s$}
{$8,5\ m/s$}
{\True $4,77\ m/s$}
{$5,17\ m/s$}
\loigiai{

}
\end{ex}
\begin{ex}%book
Hai xe máy chuyển động thẳng cùng xuất phát từ bưu điện đến hai vị trí khác nhau. Xe thứ nhất đi được quãng đường $1\ km$ hết $1,5$ phút, xe thứ hai đi được $0,75\ km$ hết $1$ phút. Tốc độ $v_1$ của xe thứ nhất và tốc độ $v_2$ của xe thứ hai được thể hiện ở liên hệ nào dưới đây?
\choice
{$4 v_1=5 v_2$}
{$8 v_1=9 v_2$}
{\True $9 v_1=8 v_2$}
{$5 v_1=4 v_2$}
\haidong{12}
\loigiai{
}
%<MyLT1>
\end{ex}
\begin{ex}
Một vật chuyển động với quãng đường vật đi được $s=13,8 \pm 0,2\ m$ trong khoảng thời gian $t=4,0 \pm 0,3\ s$. Phép đo tốc độ trung bình có sai số tỉ đối bằng
\choice
{$3,5\%$}
{$6,5\%$}
{\True $8,9\%$}
{$2,2\%$}
\haidong{13}
\loigiai{
$v=\dfrac{d}{t} \Rightarrow \delta(v)=\dfrac{\Delta v}{\overline{v}} 100\%=\left(\dfrac{\Delta d}{\overline{d}}+\dfrac{\Delta t}{\overline t}\right) 100\%=\left(\dfrac{0,2}{13,8}+\dfrac{0,3}{4,0}\right) 100\%=\left(\dfrac{0,8+4,14}{13,8.4,0}\right) 100\%=8,9\%$
}
\end{ex}
\begin{ex}%book
Trong bộ thí nghiệm đo tốc độ chuyển động của viên bi thép, cổng quang điện có vai trò giống như bộ phận nào?
\choice
{Công tắc bấm thả viên bi}
{Đồng hồ đo điện số}
{Công thắc điều khiển mở}
{\True Công tắc điều khiển đóng/mở đồng hồ đo}

\loigiai{
}
\haidong{7} \end{ex}
\begin{ex}%book
khi sử dụng đồng hồ đo thời gian hiện số và cổng quang điện có
\choice
{\True Kết quả đo chính xác, giảm thiểu sai số}
{Kết quả đo chưa chính xác, sai số nhiều}
{Đo chưa chính xác, thiết bị cồng kềnh}
{Kết quả đo chính xác, thiết bị nhỏ gọn}

\loigiai{
}
\haidong{7} \end{ex}
\begin{ex}%book
Chức năng của động hồ đo thời gian hiện số.

\choice
{\True $MODE$ $\mathbf{B}$: Đo thời gian vật chắn cổng quang điện nối với ổ $B$}
{$MODE$ $\mathbf{A} \leftrightarrow \mathbf{B}$: Đo tổng của hai khoảng thời gian vật chắn cổng quang điện nối với ổ A và vật chắn cổng quang điện nối với ổ $B$}
{$MODE$ $\mathbf{A}$: Đo thời gian từ lúc vật chuyển động đến khi vật chắn cổng quang điện nối với ổ $A$}
{$MODE$ $\mathbf{A} + \mathbf{B}$: Đo thời gian vật chuyển động từ cổng quang điện nối với ổ $A$ tới cổng quang điện nối với ổ $B$}
\loigiai{
}
\haidong{7} \end{ex}
\begin{ex}%book
Trong thí nghiệm đo tốc độ chuyển động của viên bi thép, đường kính viên bi được đo bằng dụng cụ nào sau đây?
\choice
{\includegraphics[width=3cm]{images/6-1}}
{\True 
\includegraphics[width=5cm]{images/3_4}
}
{		\includegraphics[width=4cm]{images/3_5}}
{	\includegraphics[width=5cm]{images/3_6}}

\loigiai{	}
\haidong{7} %<MyLT1>
\end{ex}
\begin{ex}%book
Để đo tốc độ tức thời của viên bi chuyển động thẳng, người ta cho viên bi chuyển động qua cổng quang điện trong khoảng thời gian t rất nhỏ. Độ dịch chuyển của bi trong khoảng thời gian t là d. Kết quả đo d và t như sau: $d=10,08 \pm 0,02\ mm$ và $t=0,036 \pm 0,001\ s$. Sai số tương đối (tỉ đối) của phép đo vận tốc của viên bi là
\choice
{$\delta v=7,1\%$}
{$\delta v=0,2\%$}
{\True $\delta v=3,0\%$}
{$\delta v=2,6\%$}
\haidong{13}
\loigiai{
}
\haidong{7} \end{ex}























































\subsubsection*{PHẦN 2. Thí sinh trả lời từ câu 1 đến câu 4. Trong mỗi ý a), b), c), d) ở mỗi câu, thí sinh chọn đúng hoặc sai.}
\setcounter{ex}{0}
\begin{ex}%book
Một người bơi dọc theo chiều dài $50\ m$ của bể bơi hết $20\ s$, rồi quay về lại chỗ xuất phát trong $22\ s$.
\choiceTF[t]
{\True Trong lần bơi đầu, vận tốc trung bình là $2,5\ m/s$}
{Trong lần bơi về, vận tốc trung bình là $2,38\ m/s$}
{\True Trong toàn bộ quá trình đi và về, độ dịch chuyển bằng $0$}
{Trong toàn bộ quá trình đi và về, vận tốc trung bình là $2,38\ m/s$}
\loigiai{ 
\begin{itemchoice}
\itemch
\itemch
\itemch
\itemch
\end{itemchoice}
}
\end{ex}
\begin{ex}%Tailieuchuan
Cho ba vật chuyển động đều. Vật thứ nhất đi được quãng đường $27\ km$ trong $30$ phút, vật thứ hai đi quãng đường $48\ m$ trong $3$ giây, vật thứ ba đi với tốc độ $60\ km/h$.
\choiceTF[t]
{\True Vật thứ nhất đi với tốc độ trung bình $54 \ km/h$} 
{\True Vật thứ hai đi với tốc độ trung bình $16 \ m/s$} 
{Vật thứ hai chuyển động chậm nhất, vật thứ ba chuyển động nhanh nhất} 
{\True Nếu đi cùng một quãng đường thì vật thứ ba chuyển động ít thời gian nhất} 
\loigiai{ 
\begin{itemchoice}
\itemch
\itemch
\itemch
\itemch
\end{itemchoice}
}
\end{ex}
\begin{ex}%Tailieuchuan
Hai xe ô tô khởi hành cùng một lúc từ hai địa điểm $A$ và $B$, cùng chuyển động về địa điểm $C$. Biết xe $(1)$ khởi hành từ $A$ đi với tốc độ $60 \ km/h$, xe $(2)$ từ $B$ chuyển động với tốc độ $55 \ km/h$ và quãng đường $AC$ dài $120 \ km$, $BC$ dài $80\ km$.
\choiceTF[t]
{\True Xe $(1)$ chuyển động đến $C$ trong thời gian $2$ giờ} 
{\True Xe $(2)$ đến $C$ trước xe $A$} 
{Để hai xe đến cùng nhau thì xe $(2)$ phải chuyển động với tốc độ $50 \ km/h$} 
{Nếu xe $(1)$ và xe $(2)$ cùng khởi hành tại $B$ với tốc độ như ban đầu, thì xe $A$ sẽ đến $C$ trước xe $B$ $10$ phút} 
\loigiai{ 
\begin{itemchoice}
\itemch
\itemch
\itemch
\itemch
\end{itemchoice}
}
\end{ex}
\begin{ex}%Tailieuchuan
\immini[thm]{Một người đi bộ đi với tốc độ không đổi dọc theo nửa đường tròn có bán kính $5\ m$, từ $A$ đến $B$ như hình vẽ với thời gian đi là $6,0 \ s$.}{
\begin{tikzpicture}[scale=1,thick,>=stealth',draw=Mapcolor]
\path
(4.5,-0.5)coordinate(A)
;
\draw[<-,Mapcolor] (A) arc (0:180:1 and 1);
\node at (2.5,-1) {$A$};
\node at (4.5,-1) {$B$};

\end{tikzpicture}}
\choiceTF[t]
{Quãng đường người đó đã đi chính là chu vi đường tròn} 
{Độ dịch chuyển của người đó bằng quãng đường người đó đã đi} 
{Tốc độ trung bình của người đó là $\dfrac{5\pi}{12}\ m/s$} 
{Vận tốc trung bình của người đó là $\dfrac{5}{6} \ m/s$} 
\loigiai{ 
\begin{itemchoice}
\itemch
\itemch
\itemch
\itemch
\end{itemchoice}
}
\end{ex}
\subsubsection*{PHẦN 3. Thí sinh trả lời từ câu 1 đến câu 6.}
\setcounter{ex}{0}
\begin{ex}%Tailieuchuan
Một ô tô chuyển động thẳng đều. Biết rằng ô tô chuyển động theo chiều âm với vận tốc có độ lớn $36\ km / h$ và ở thời điểm $1,5\ h$ thì vật có tọa độ $6\ km$. Vị trí của vật sau $1\ h$ là bao nhiêu $km$?
\shortans[oly]{24}
\loigiai{
}
\end{ex}
\begin{ex}%Tailieuchuan
Một tàu ngầm sử dụng hệ thống phát sóng âm để đo độ sâu của biển. Hệ thống phát ra các sóng âm và đo thời gian quay trở lại của sóng âm sau khi chúng bị phản xạ tại đáy biển. Tại một vị trí trên mặt biển, thời gian mà hệ thống ghi nhận được là $0,13\ s$ kể từ khi sóng âm được truyền đi. Biết tốc độ truyền sóng âm trong nước khoảng $1500\ m/s$. Nước biển sâu bao nhiêu mét? 
\shortans[oly]{97,5}
\loigiai{
}
\end{ex}
\begin{ex}%Tailieuchuan
Một ô tô xuất phát từ $A$ lúc $8$ giờ sáng chuyển động thẳng đều tới $B$ lúc $10$ giờ $30$ phút, đến $C$ lúc $12$ giờ $30$ phút. Khoảng cách từ $A$ đến $B$ là $175\ km$. Khoảng cách từ $B$ đến $C$ bao nhiêu km?
\shortans[oly]{140}
\loigiai{
}
\end{ex}
\begin{ex}
Một người bắt đầu cho xe máy chạy trên một đoạn đường thẳng. Trong $10$ giây đầu, xe chạy được quãng đường $50\ m$, trong $10$ giây tiếp theo xe chạy được quãng đường $100\ m$. Tốc độ trung bình của xe máy trong $20$ giây đầu tiên là bao nhiêu $km/h$?
\shortans[oly]{27}
\loigiai{
Tốc độ trung bình của xe máy trong 20 giây đầu tiên: $v_{tb}=\dfrac{s}{t}=\dfrac{50+100}{10+10}=7,5 ~m/s=27 ~km / h$
}
\end{ex}
\begin{ex}
Cho một xe ô tô chạy trên một quãng đường trong $5$ giờ. Biết $2$ giờ đầu xe chạy với tốc độ trung bình $60\ km/h$ và $3$ giờ sau xe chạy với tốc độ trung bình $40\ km/h$. Tốc trung bình của xe trong suốt thời gian chuyển động là bao nhiêu $km/h$?
\shortans[oly]{48}
\loigiai{		
\begin{itemize}
\item Quãng đường đi trong 2 h đầu: $S_1=v_1 t_1=2.60=120 ~km$
\item Quãng đường đi trong 3 h sau: $S_2=v_2 t_2=3.40=120 ~km$
\item Tốc độ trung bình của xe trong suốt thời gian chuyển động: $v=\dfrac{s_1+s_2}{t}=\dfrac{120+120}{5}=48 ~km / h$.
\end{itemize}
}
\end{ex}
\begin{ex}
Một ô tô chuyển động trên đường thẳng. Tại thời điểm $t_1$, ô tô ở cách vị trí xuất phát $5\ km$. Tại thời điểm $t_2$, ô tô ở cách vị trí xuất phát $12\ km$. Từ $t_1$ đến $t_2$, độ dịch chuyển của ô tô đã thay đổi một đoạn bằng bao nhiêu $km$?
\shortans[oly]{7}
\loigiai{
Độ dịch chuyển của ô tô là: $d=12-5=7\ km$
}
\end{ex}

